\subsection{Operational Saturation Scale}
  \label{subsec:operational_saturation_scale}

  The bounded-response behavior introduced above can be associated with
  an operational saturation scale.
  In the underlying relational description, the invariant bound $b_{\chi}$
  limits the maximal rate at which admissible configurations can be updated.
  In the effective macroscopic regime, this constraint manifests as a
  characteristic acceleration proxy $a_\star$ marking the onset of
  non-linear response.

  Dimensional consistency requires that this threshold be inversely
  proportional to the descriptive resolution $\ell(z)$ of the effective
  state $U$.
  We therefore parametrize
  \begin{equation}
    \label{eq:astar_param}
    a_\star(z) \sim \kappa \frac{b_{\chi}}{\ell(z)},
  \end{equation}
  where $\kappa$ is a dimensionless normalization factor encoding the
  convention adopted for the projective mapping.
  In the late-time low-density regime considered in this section,
  $a_\star$ may be treated as approximately constant.
  Its possible redshift dependence becomes relevant only in the
  primordial regime and will be discussed in Section~\ref{sec:predictions}.
